\documentclass[11pt,a4paper]{article}
\usepackage[T1]{fontenc}
\usepackage[utf8]{inputenc}
\usepackage{lmodern,amsmath,amssymb}
\usepackage[margin=25mm]{geometry}
\usepackage[hidelinks]{hyperref}
\hypersetup{pdftitle={Two exact identities for the magnetic energy, a ground-state sum rule, and why pointwise G},pdfauthor={navstokgap research working notes}}
\newcommand{\tightlist}{\setlength{\itemsep}{0pt}\setlength{\parskip}{0pt}}
\setlength{\emergencystretch}{3em}
\setcounter{secnumdepth}{0}
\begin{document}
\hypertarget{two-exact-identities-for-the-magnetic-energy-a-ground-state-sum-rule-and-why-pointwise-gibbs-domination-fails-as-agmon-does}{%
\section{Two exact identities for the magnetic energy, a ground-state
sum rule, and why pointwise Gibbs domination fails as Agmon
does}\label{two-exact-identities-for-the-magnetic-energy-a-ground-state-sum-rule-and-why-pointwise-gibbs-domination-fails-as-agmon-does}}

The magnetic energy of the Kogut--Susskind Hamiltonian satisfies two
exact relations on the configuration space \(G^{\mathcal E}\):
\[\Delta V=4C_2\big(N|\mathcal P|-V\big),
\qquad
\big|\nabla V\big|^2\ \le\ 16\,V\quad(SU(2)),\] the first because
\(V-N|\mathcal P|\) is an eigenfunction of the Laplace--Beltrami
operator with eigenvalue \(-4C_2\), the second from the completeness
relation together with \(\sin^2\theta=(1-\cos\theta)(1+\cos\theta)\).
They give an exact sum rule for the ground state,
\[A'\!\int\!\big(V-N|\mathcal P|\big)\big|\nabla\Omega\big|^2d\mu
+A'\big(2C_2-K\big)\big\langle V-N|\mathcal P|\big\rangle
+B'\operatorname{Var}_\Omega(V)=0,
\qquad K=\!\int\!|\nabla\Omega|^2,\] and they make the
comparison-principle attempt at a pointwise bound explicit: with
\(\varphi=e^{-\lambda V/2}\) the operator \(L=-A'\Delta+(B'V-e_0)\)
satisfies \(L\varphi\ge0\) precisely above a threshold
\(V_*\simeq N|\mathcal P|\), so the maximum principle bounds \(\Omega\)
only above the extensive Haar mean. \textbf{The pointwise Gibbs
domination therefore fails for the same reason as the Agmon bound}: the
ground-state energy \(e_0\) is extensive and enters every comparison, so
a local excess is never in the region where the comparison bites. The
obstruction is that the Hamiltonian ground state is characterized by a
global variational principle while the Euclidean weight \(e^{-S_w}\) is
local by construction, and no manipulation of the eigenvalue equation
alone repairs that. Constants explicit; the two identities and the sum
rule are exact; nothing promoted.

\hypertarget{the-laplacian-of-the-magnetic-energy}{%
\subsection{1. The Laplacian of the magnetic
energy}\label{the-laplacian-of-the-magnetic-energy}}

\textbf{Proposition 1.} On \(\mathcal M=G^{\mathcal E}\) with the
product bi-invariant metric,
\[\Delta V=4C_2(R_{\rm f})\big(N|\mathcal P|-V\big),\] where
\(C_2(R_{\rm f})\) is the quadratic Casimir of the defining
representation. Equivalently \(W=V-N|\mathcal P|\) satisfies
\(\Delta W=-4C_2W\): the magnetic energy, shifted by its Haar mean, is
an eigenfunction of the Laplacian.

\emph{Proof.} Fix a plaquette \(p\) and one of its four links \(\ell\).
Writing \(\operatorname{tr}U_p=\sum_{ab}(U_\ell)_{ab}M_{ba}\) with \(M\)
the ordered product of the other three links, each matrix element of
\(U_\ell\) in the defining representation is an eigenfunction of
\(\Delta_\ell\) with eigenvalue \(-C_2(R_{\rm f})\), by Peter--Weyl.
Hence \(\Delta_\ell\operatorname{tr}U_p=-C_2\operatorname{tr}U_p\), and
summing over the four links of \(p\),
\(\Delta\operatorname{Re}\operatorname{tr}U_p=-4C_2\operatorname{Re}\operatorname{tr}U_p\).
Summing over plaquettes,
\(\Delta V=\Delta\sum_p(N-\operatorname{Re}\operatorname{tr}U_p) =4C_2\sum_p\operatorname{Re}\operatorname{tr}U_p=4C_2(N|\mathcal P|-V)\).
\(\square\)

\hypertarget{the-gradient-of-the-magnetic-energy}{%
\subsection{2. The gradient of the magnetic
energy}\label{the-gradient-of-the-magnetic-energy}}

\textbf{Proposition 2.} For \(SU(2)\), \(|\nabla V|^2\le16\,V\)
pointwise on \(\mathcal M\). For general \(SU(N)\),
\(|\nabla V|^2\le16\,N\,V\).

\emph{Proof.} For one link of one plaquette, the completeness relation
of \href{lattice-gap-upper-bounds.md}{the upper-bound note} Corollary 2
gives \[\sum_a\big|X^\ell_a\operatorname{tr}U_p\big|^2
=\tfrac12\Big[\operatorname{tr}(MM^\dagger)-\tfrac1N\big|\operatorname{tr}M\big|^2\Big]
=\tfrac12\Big[N-\tfrac1N\big|\operatorname{tr}U_p\big|^2\Big].\] For
\(SU(2)\), writing \(\operatorname{tr}U_p=2\cos\theta\), this is
\(1-\cos^2\theta=\sin^2\theta=(1-\cos\theta)(1+\cos\theta)\), while
\(V_p=2-2\cos\theta\), so
\[\big|\nabla_\ell\operatorname{Re}\operatorname{tr}U_p\big|^2
\le\sin^2\theta=\frac{V_p}2\Big(2-\frac{V_p}2\Big)\le V_p .\] Each link
lies in four plaquettes, so by Cauchy--Schwarz
\(|\nabla_\ell V|^2\le4\sum_{p\ni\ell}V_p\), and summing over links,
with four links per plaquette,
\(|\nabla V|^2\le4\sum_\ell\sum_{p\ni\ell}V_p=16\sum_pV_p=16V\). The
general case replaces the bound \(\sin^2\theta\le V_p\) by
\(\tfrac12[N-|\operatorname{tr}U_p|^2/N]\le NV_p\), which follows from
\(|\operatorname{tr}U_p|\ge N-V_p\) when \(V_p\le N\) and from the
trivial bound otherwise. \(\square\)

Both identities vanish at \(U_p=1\): a configuration with no magnetic
energy has no magnetic gradient, as it must.

\hypertarget{an-exact-ground-state-sum-rule}{%
\subsection{3. An exact ground-state sum
rule}\label{an-exact-ground-state-sum-rule}}

\textbf{Proposition 3.} With \(W=V-N|\mathcal P|\),
\(K=\int|\nabla\Omega|^2d\mu\),
\(\langle\cdot\rangle=\langle\Omega,\cdot\,\Omega\rangle\) and
\(\operatorname{Var}_\Omega(V)=\langle V^2\rangle-\langle V\rangle^2\),
\[A'\!\int W\big|\nabla\Omega\big|^2d\mu
+A'\big(2C_2-K\big)\langle W\rangle
+B'\operatorname{Var}_\Omega(V)=0 .\]

\emph{Proof.} Multiply \(-A'\Delta\Omega+(B'V-e_0)\Omega=0\) by
\(W\Omega\) and integrate. Since
\[\int W\Omega\,\Delta\Omega=-\int W|\nabla\Omega|^2-\int\Omega\,\nabla W\!\cdot\!\nabla\Omega,\]
\[\int\Omega\,\nabla W\!\cdot\!\nabla\Omega=-\tfrac12\int(\Delta W)\Omega^2=2C_2\langle W\rangle,\]
by Proposition 1, so
\(\int W\Omega\Delta\Omega=-\int W|\nabla\Omega|^2-2C_2\langle W\rangle\)
and the equation becomes
\(A'\int W|\nabla\Omega|^2+2A'C_2\langle W\rangle+B'\langle VW\rangle-e_0\langle W\rangle=0\).
Now
\(e_0=A'K+B'\langle V\rangle=A'K+B'(\langle W\rangle+N|\mathcal P|)\)
and
\(\langle VW\rangle=\langle W^2\rangle+N|\mathcal P|\langle W\rangle\),
so the last two terms combine into
\(-A'K\langle W\rangle-B'\langle W\rangle^2+B'\langle W^2\rangle\).
\(\square\)

The sum rule ties the magnetic variance of the ground state to its
kinetic energy and to the correlation between magnetic energy and
kinetic density. At \(g\to\infty\) every term vanishes with \(B'\), and
at finite coupling it is one exact constraint on the ground-state
measure.

\hypertarget{pointwise-gibbs-domination-and-where-it-stops}{%
\subsection{4. Pointwise Gibbs domination and where it
stops}\label{pointwise-gibbs-domination-and-where-it-stops}}

Try \(\varphi=e^{-\lambda V/2}\), \(\lambda>0\), as a comparison
function for \(L=-A'\Delta+(B'V-e_0)\). Using
\(\Delta\varphi=\varphi[\tfrac{\lambda^2}4|\nabla V|^2-\tfrac\lambda2\Delta V]\)
and Propositions 1--2, \[\frac{L\varphi}{\varphi}
=B'V-e_0-A'\frac{\lambda^2}4|\nabla V|^2+A'\frac\lambda2\Delta V
\ \ge\ \Big[B'-4A'\lambda^2-2A'\lambda C_2\Big]V
+\Big[2A'\lambda C_2N|\mathcal P|-e_0\Big],\] for \(SU(2)\), using
\(|\nabla V|^2\le16V\) and Proposition 1.

\textbf{Proposition 4.} If \(\lambda\) satisfies
\(B'>4A'\lambda^2+2A'\lambda C_2\), then \(L\varphi\ge0\) on
\(\{V\ge V_*\}\) with
\[V_*=\frac{\big(e_0-2A'\lambda C_2N|\mathcal P|\big)_+}{B'-4A'\lambda^2-2A'\lambda C_2},\]
and the maximum principle gives
\[\Omega(U)\ \le\ \Big(\max_{\{V<V_*\}}\frac{\Omega}{\varphi}\Big)\,e^{-\lambda V(U)/2}
\qquad\text{for all }U .\]

\emph{Proof.} Set \(u=\Omega/\varphi\). From \(L\Omega=0\),
\(0=-A'\varphi\Delta u-2A'\nabla u\!\cdot\!\nabla\varphi+u\,L\varphi\).
At an interior maximum of \(u\) one has \(\nabla u=0\) and
\(\Delta u\le0\), so the first two terms are \(\ge0\) and
\(u\,L\varphi\le0\); since \(u>0\), the maximum lies in
\(\{L\varphi\le0\}\subset\{V<V_*\}\). \(\square\)

\textbf{Where it stops.} With \(e_0\le B'N|\mathcal P|\), the threshold
obeys
\[V_*\ \le\ N|\mathcal P|\;\frac{B'-2A'\lambda C_2}{B'-4A'\lambda^2-2A'\lambda C_2}
\ \xrightarrow[\lambda\to0]{}\ N|\mathcal P| ,\] so the comparison
function controls \(\Omega\) only relative to its values on the set
where \(V\) is below the \textbf{extensive} Haar mean. The bound that
results, \(\Omega\le C\,e^{-\lambda(V-V_*)/2}\) with \(C\) determined on
\(\{V<V_*\}\), is again a statement about the global excess. A local
excess of \(n\) plaquettes leaves \(V\) far below
\(V_*\simeq N|\mathcal P|\), so the bound says nothing about it, exactly
as in \href{agmon-global-not-local.md}{the global/local note}.

\hypertarget{the-structural-reason-stated-once}{%
\subsection{5. The structural reason, stated
once}\label{the-structural-reason-stated-once}}

Both attempts, Agmon and comparison, use the eigenvalue equation
\(-A'\Delta\Omega+(B'V-e_0)\Omega=0\), in which \(e_0\) is the
\textbf{total} ground-state energy and is extensive. Any inequality
derived from it compares the total potential with the total energy, so
it distinguishes configurations only by their global excess. The
Euclidean weight \(e^{-S_w}=\prod_pe^{-(2/g_E^2)V_p}\) factorizes over
plaquettes and therefore separates configurations locally by
construction, with no reference to a total energy.

This is the precise content of the observation in
\href{typical-field-strength-window.md}{the typical-field note}: the
Hamiltonian formulation has a state and a global eigenvalue equation
where the Euclidean formulation has a local weight. The three notes
since then have tried the two standard ways of extracting a local
statement from a global equation, and both stop at the same place. A
local statement would need an input of a different kind, for instance a
proof that the ground state is a Gibbs state for a local potential, or
the transfer-matrix identification of \(\Omega^2\) with a Euclidean
measure on a time slice, which is the route the constructive programme
takes.

\hypertarget{consequence-for-state}{%
\subsection{6. Consequence for STATE}\label{consequence-for-state}}

The question named in \href{agmon-global-not-local.md}{the global/local
note} §5 is answered in the negative by Proposition 4: pointwise Gibbs
domination of the Kogut--Susskind ground state does not follow from the
eigenvalue equation, for the same extensivity reason as the Agmon bound.
What the attempt produced instead is worth keeping: the exact relations
\(\Delta V=4C_2(N|\mathcal P|-V)\) and \(|\nabla V|^2\le16V\), and the
exact ground-state sum rule of Proposition 3, which constrains the
magnetic variance in terms of kinetic quantities. The line of attack
that remains is the transfer-matrix identification of the ground-state
measure with a Euclidean measure, which leaves the Hamiltonian framework
by design.
\end{document}
